% Fragmento público generado desde propietarios íntegros.
% Residencia: 52.10.2.
% Fuente: ../incoming/radion_monografia_20260827/manuscrito/sections/11_modularidad_orbitas.tex:85-187
\subsubsection{Conservation surfacique dans les orbites de Kepler et de Sommerfeld}

\paragraph{Équivalence surfacique des réalisations orbitale et oscillatoire}

L'ellipse de Kepler vit dans l'espace de configuration et répond à un potentiel
central. L'ellipse du radion vit dans l'espace des phases et répond à une forme
d'action. Ce sont des objets différents. Elles partagent une loi surfacique parce que les deux
réalisations possèdent une deux-forme conservée.

Pour une orbite képlérienne,
\begin{equation}
r(\phi)=\frac{p}{1+e\cos(\phi-\phi_0)},
\qquad
\frac{\diff A}{\diff t}=\frac{\ell}{2m}.
\label{eq:kepler}
\end{equation}
La troisième loi prend la forme
\[
T^2=\frac{4\pi^2\mu}{K}a^3.
\]
Pour l'oscillateur en phase,
\begin{equation}
J=\oint p\,\diff q=\frac{2\pi E}{\omega}.
\label{eq:osc-action}
\end{equation}
Si la cellule élémentaire est \(J=\hfull\), alors
\begin{equation}
E=\hret\omega,\qquad \frac{J}{2\pi}=\hret.
\label{eq:E-hbarw}
\end{equation}
Telle est la convergence nette entre action de Planck, fréquence angulaire et
radion.

Sommerfeld quantifie les intégrales d'action sur les cycles. La lecture HMT ajoute
la mémoire du cycle, du feuillet et du retour qui survivent. Dans un ruban de Möbius,
un lacet visible peut porter \(\hfull/2\) et deux lacets restaurer \(\hfull\).
Ce facteur topologique ne s'ajoute pas à l'indice de Maslov d'une quantification EBK ;
ils relèvent de corrections ayant des propriétaires distincts.

\paragraph{Holonomie orbitale et propagations avancée et retardée}

Une orbite qui accumule de l'holonomie satisfait une condition du type
\begin{equation}
\epsilon_\gamma\exp\left(\frac{iJ_\gamma}{\hret}\right)=1.
\label{eq:holonomy-action}
\end{equation}
Le signe \(\epsilon_\gamma\) conserve le feuillet ; \(J_\gamma\) conserve
l'action. L'avance ou le retard apsidal s'interprète comme découplage entre
la clôture angulaire visible et le retour complet. Le radion fournit
l'unité orientée avec laquelle cette différence est mesurée.

\subsubsection{Composition des amplitudes de Schrödinger et somme des chemins}

\paragraph{Quatre réalisations d'une même phase}

Une fois \(\hret\) fixé, différentes équations conventionnelles reconnaissent
des aspects du même transport :
\begin{enumerate}
\item Schrödinger publie la phase au moyen de
\(i\hret\partial_t\psi=H\psi\).
\item Pauli ajoute le double feuillet de spin et son couplage magnétique.
\item Dirac linéarise la relation énergie--quantité de mouvement et organise particule et
antiparticule dans une représentation causale.
\item La somme des chemins attribue à chaque histoire la phase
\(\exp(iS[\gamma]/\hret)\).
\end{enumerate}
Ces théories ne génèrent pas rétrospectivement \(\hret\). Elles reçoivent l'unité
d'action et réalisent différents contenus de l'état enrichi.

\paragraph{Fonctionnelle de phase sur l'espace des chemins}

Dans la formulation de Feynman,
\begin{equation}
\mathcal K(b,a)=\int_{\gamma:a\to b}
\exp\left(\frac{i}{\hret}S[\gamma]\right)\mathcal D\gamma.
\label{eq:path-integral}
\end{equation}
Un incrément d'action \(\hret\) change la phase d'un radian. Un incrément
\(\hfull\) la change d'un tour. Le théorème
\(\mathcal A(\theta)=\hret\theta\) transforme cette relation en géométrie
d'aire : la phase d'un chemin peut se lire comme action balayée en unités de
radion.

L'inversion \(C_M=-\operatorname{rev}\) sur les mots dodécaphasiques et
l'identité \eqref{eq:time-reversal} organisent des chemins conjugués. Dans une
réalisation d'absorbeur, on sépare
\[
G_{\mathrm{sym}}=\tfrac12(G_{\mathrm{ret}}+G_{\mathrm{adv}}),
\qquad
G_{\mathrm{odd}}=G_{\mathrm{ret}}-G_{\mathrm{adv}}.
\]
Le premier secteur est pair sous l'échange retardé--avancé ; le second est
impair. Cela reproduit la séparation entre énergie/forme paire et action
orientée impaire du radion lorsque le transducteur préserve la frontière et la forme
symplectique.

\begin{statusbox}
La théorie classique de l'absorbeur et l'interprétation de Feynman--Stückelberg des
antiparticules sont des réalisations historiques différentes. La première utilise une
interaction temporellement symétrique ; la seconde réorganise la propagation dans la
théorie relativiste. HMT apporte un antécédent commun de réversion, mais ne les
fusionne pas sans une application de propagateurs.
\end{statusbox}
